\section{Conclusion}

We revisited RLVR from the perspective of reward distribution estimation and identified point estimation under limited rollouts as a key source of sample inefficiency and variance collapse.
To this end, we proposed Discounted Beta--Bernoulli (DBB) reward estimation for RLVR, which leverages historical reward information for accurate estimation.
Although biased, the DBB estimator achieves substantially lower variance and, under mild non-stationarity of rewards, a lower mean squared error than point estimation; it also theoretically avoids variance collapse and preserves informative training signals without additional rollouts or replay.
Experiments across two model scales demonstrate consistent improvements over baselines on all in-distribution and out-of-distribution benchmarks.
Our findings highlight reward distribution estimation as a critical yet underexplored component of effective RLVR and suggest that principled estimator design offers a promising direction for improving large-scale reinforcement learning of language models.
As future work, we plan to design a dynamic discount factor $\lambda$ that adapts across different stages of training.
We also view DBB as one instance of a broader discounted Bayesian reward-estimation framework.
A natural next step is to pair the discounted update with other likelihoods (e.g., a discounted Gaussian) to handle continuous or dense process rewards.


\section*{Impact Statements}
This paper presents work whose goal is to advance the field of machine learning. There are many potential societal consequences of our work, none of which we feel must be specifically highlighted here.

\section*{Acknowledgements}
This work was supported by the Institute of Information \& Communications Technology Planning \& Evaluation (IITP) grant funded by the Korea government (MSIT) (RS-2019-II190075, Artificial Intelligence Graduate School Program(KAIST)) and National Research Foundation of Korea (NRF) grant funded by the Korea government (MSIT)
(RS-2023-00209060, A Study on Optimization and Network Interpretation Method for LargeScale Machine Learning).